\documentclass[a4paper]{article}
\usepackage[pdftex]{graphicx}
\usepackage[utf8]{inputenc}
\usepackage{enumerate}
\usepackage{siunitx}
\sisetup{locale=DE} 
\usepackage{href-ul}
\hypersetup{
	colorlinks=true,
	linkcolor=blue,
	urlcolor=blue}
\usepackage{icomma}
\usepackage{amssymb}
\hyphenation{Tem-pe-ra-tur-ska-la}
\usepackage{geometry}
\geometry{a4paper, top=15mm, left=15mm, right=15mm, bottom=15mm,
	headsep=10mm, footskip=12mm}
%Source: img_0112

% Start the document
\begin{document}
%Titel
{\bf Schulaufgabe aus der Physik }\\
\begin{enumerate}[1.]
%Aufgabe 1
{\bf \item Erläutere, was man unter folgenden Begriffen versteht:}

a) Wärme \hspace{1 cm} b) Innere Energie  \hspace{1 cm} c) Anomalie des Wassers 

%Aufgabe 2
{\bf \item Nenne drei verschiedene Auswirkungen, die die Veränderung der inneren Energie haben kann.}

%Aufgabe 3
{\bf \item Warum ist es möglich, mit dem Verfahren von Anders Celsius eine Temperaturskala festzulegen?}

%Aufgabe 4
{\bf \item Beschreibe den Vorgang der Wärmeleitung im Teilchenmodell.}

%Aufgabe 5
{\bf \item Beschreibe ausführlich die Entstehung der Wärmeströmung der Luft in einem Raum, der mit Hilfe eines Heizkörpers erwärmt wird.}

%Aufgabe 6
{\bf \item Tee in einem Thermosgefäß bleibt lange heiß.} \\
Beschreibe, durch welche physikalische Maßnahmen dies erreicht wird, und erkläre kurz, warum das durch diese Maßnahme jeweils möglich ist.
 
%Aufgabe 7
\begin{minipage}{0.65\textwidth}
{\bf \item In einem Versuch wird der Zusammenhang zwischen der einem Körper zugeführten Wärmemenge $W_{th}$} und der Temperaturänderung des Körpers untersucht. Dazu wird an einem Metallzylinder mit dem Durchmesser $d=\SI{5,6}{\centi\meter}$ durch Drehen des Zylinders Reibungsarbeit $W_R$ verrichtet. Die Schnur, die um den Zylinder gewickelt ist, wird von einem Masse\-stück mit der Gewichtskraft $F_G= \SI{25}{\newton}$ ge\-spannt. n ist die Anzahl der Umdrehungen. 
\end{minipage}
\hfill
\begin{minipage}{0.3\textwidth}
\includegraphics[width=4cm]{rolle112.pdf}
\end{minipage}

\begin{enumerate}[a)]
\item Berechne die Wärmemenge, die bei einer Umdrehung dem Zylinder zugeführt wird. \\

Es ergibt sich folgende Messtabelle:

\renewcommand{\arraystretch}{2}
\begin{tabular}{ccccc}
$\Delta \vartheta ~\textnormal{in}\SI{}{\celsius}$ & 0,51 & 0,98 & 1,5 & 2,3 \\
\hline
n & 100 & 200 & 300 & 450 \\
\hline
$W_{th}  \textnormal{in}~\SI{}{\kilo\joule}$ & & & & 
\end{tabular}

\item Ergänze die Tabelle durch $W_{th} ~\textnormal{in}\SI{}{\kilo\joule}$ und werte sie zeichnerisch aus.\\
$[\textnormal{y-Achse}~= W_{th}; \SI{1}{\centi\meter}~ \hat{=}~\SI{0,25}{\kilo\joule};~
\textnormal{x-Achse}~=\Delta \vartheta  ;~ \SI{1}{\centi\meter}~ \hat{=}~\SI{0,25}{\celsius}] $
Welcher Zusammenhang ergibt sich?\\

\item Wie kannst du den Zusammenhang auch rechnerisch erkennen? Welche Masse hat der Zylinder, wenn er aus Messing mit der  spezifischen Wärmekapazität
 $c=\SI{0,384}{\kilo\joule\over \kilogram\cdot\kelvin}$ ist?
\end{enumerate}

\end{enumerate}

\fbox{
	\begin{minipage}{0.5\textwidth}
		Zur Lösung bitte \href{https://www.okuyakl.de/phys/p9thcL112/ll112.pdf}{hier klicken} oder den QR-Code scannen.\\
	Weitere Arbeitsblätter gibt es unter 
	
	\href{https://www.okuyakl.de}{www.okuyakl.de}
	\end{minipage}
	\hfill
	\begin{minipage}{0.4\textwidth}
		\includegraphics[width=1.5 cm]{../../viecher/zwe03}
		\includegraphics[width=3 cm]{qr112}
		\includegraphics[width=2 cm]{../../viecher/afanticon1}
		
	\end{minipage}}

\end{document}%L Lösung ------------------------------------------------------------------------
